% =====================================================================
\chapter{Del perceptrón al perceptrón multicapa}
\label{tema:1}
% =====================================================================

Este tema presenta la unidad básica de una red neuronal, el perceptrón, su regla de aprendizaje y el teorema que garantiza su convergencia, y muestra por qué una sola unidad no basta: la función XOR no es linealmente separable. La solución, añadir una capa oculta, conduce al perceptrón multicapa. La exposición sigue la §6.1 de~\cite{gbc} y las §6.1--6.3 de~\cite{jm3}; la demostración del teorema de convergencia sigue la lección~3 de~\cite{cs4780}, basada en Novikoff~\cite{novikoff}, y el recuento de parámetros, las notas «Neural Networks Part~1» de~\cite{cs231n}.

\section{La unidad y el perceptrón}

\begin{definicion}[Unidad]\label{def:unidad}
Una \emph{unidad} (o \emph{neurona artificial}) con $n$ entradas está determinada por un vector de \emph{pesos} $w\in\R^{n}$, un \emph{sesgo} $b\in\R$ y una función $\phi\colon\R\to\R$, llamada \emph{función de activación}. Para una entrada $x\in\R^{n}$ calcula la \emph{preactivación} $z=w\cdot x+b$ y la \emph{activación} $a=\phi(z)$.
\end{definicion}

\begin{definicion}[Perceptrón]\label{def:perceptron}
El \emph{perceptrón} de Rosenblatt~\cite{rosenblatt} es la unidad con activación escalón:
\[
\hat y=\ind[\,w\cdot x+b>0\,]\in\{0,1\}.
\]
Su \emph{frontera de decisión} es el hiperplano $\{x\in\R^{n}:w\cdot x+b=0\}$.
\end{definicion}

El perceptrón es, por tanto, un clasificador lineal: asigna la clase~1 al semiespacio abierto $\{w\cdot x+b>0\}$ y la clase~0 a su complementario.

\addtocounter{definicion}{1}%

\begin{definicion}[Separabilidad lineal]
Un conjunto finito de pares $(x_i,y_i)$ con $x_i\in\R^{n}$ e $y_i\in\{0,1\}$ es \emph{linealmente separable} si existen $w\in\R^{n}$ y $b\in\R$ tales que $w\cdot x_i+b>0$ cuando $y_i=1$ y $w\cdot x_i+b<0$ cuando $y_i=0$.
\end{definicion}

\section{La regla de aprendizaje de Rosenblatt}

\begin{definicion}[Regla del perceptrón]\label{def:regla}
Sea $\eta>0$ una \emph{tasa de aprendizaje}. Se parte de $w=0$, $b=0$ y se recorren los ejemplos de forma cíclica. Para cada ejemplo $(x_i,y_i)$ se calcula $\hat y_i$ y se actualiza
\[
w\leftarrow w+\eta\,(y_i-\hat y_i)\,x_i,\qquad b\leftarrow b+\eta\,(y_i-\hat y_i).
\]
El algoritmo termina cuando una pasada completa no produce ningún error.
\end{definicion}

Como $y_i-\hat y_i\in\{-1,0,1\}$, los pesos solo cambian cuando el ejemplo está mal clasificado, y el cambio desplaza la preactivación en la dirección correcta: si $y_i=1$ y $\hat y_i=0$, la nueva preactivación de $x_i$ aumenta en $\eta(\norm{x_i}^{2}+1)>0$.

\addtocounter{definicion}{1}%

Para enunciar el teorema de convergencia conviene absorber el sesgo y usar etiquetas simétricas. Sea $\tilde x_i=(x_i,1)\in\R^{n+1}$, $\theta=(w,b)$ y $s_i=2y_i-1\in\{-1,+1\}$. El ejemplo $i$ está bien clasificado si $s_i\,\theta\cdot\tilde x_i>0$, y la regla de la \cref{def:regla} se escribe $\theta\leftarrow\theta+\eta\,s_i\,\tilde x_i$ cada vez que se comete un error. En particular, toda actualización se produce en un ejemplo con $s_i\,\theta\cdot\tilde x_i\le 0$.

\begin{teorema}[Convergencia del perceptrón; Novikoff]\label{teo:novikoff}
Supóngase que existen $u\in\R^{n+1}$ con $\norm{u}=1$ y $\gamma>0$ tales que $s_i\,u\cdot\tilde x_i\ge\gamma$ para todo $i$, y sea $R=\max_i\norm{\tilde x_i}$. Entonces, cualquiera que sea el orden en que se presentan los ejemplos, el número total $M$ de actualizaciones de la regla del perceptrón cumple
\[
M\le\frac{R^{2}}{\gamma^{2}}.
\]
\end{teorema}

\begin{observacion}
La cota no depende de $\eta$ ni del número de ejemplos, sino solo de la geometría de los datos: el cociente $R/\gamma$ mide lo «estrecho» que es el pasillo que separa las dos clases. Si los datos no son linealmente separables, no existe tal $u$ y la regla no termina nunca.
\end{observacion}

\section{El problema XOR}

La función XOR vale~1 en $(0,1)$ y $(1,0)$ y vale~0 en $(0,0)$ y $(1,1)$.

\begin{proposicion}\label{prop:xor}
El conjunto $\{((0,0),0),((0,1),1),((1,0),1),((1,1),0)\}$ no es linealmente separable. En consecuencia, ningún perceptrón calcula XOR.
\end{proposicion}

Minsky y Papert~\cite{minsky} analizaron en 1969 esta y otras limitaciones de los perceptrones de una capa, y su libro contribuyó a frenar la investigación en redes neuronales durante más de una década~\cite[cap.~1]{gbc}.

\section{Una capa oculta resuelve XOR}

La salida consiste en componer dos capas. Sea $\ReLU(t)=\max(0,t)$ y considérese la red de la \cref{fig:xor}, tomada de~\cite[§6.1]{gbc}:
\begin{equation}\label{eq:xor}
h_1=\ReLU(x_1+x_2),\qquad h_2=\ReLU(x_1+x_2-1),\qquad \hat y=h_1-2h_2.
\end{equation}

\addtocounter{figure}{1}%

En el espacio de la capa oculta, las cuatro entradas se transforman en $(0,0)$, $(1,0)$, $(1,0)$ y $(2,1)$: las dos entradas con salida~1 caen en el mismo punto, y la función lineal $h_1-2h_2$ separa las clases. La capa oculta no clasifica: cambia la \emph{representación} para que una función lineal pueda hacerlo. Aprender la representación a la vez que el clasificador es la idea que distingue el aprendizaje profundo de los modelos lineales del módulo~8.

\section{El perceptrón multicapa}

\begin{definicion}[Perceptrón multicapa]\label{def:mlp}
Sean $n_0,n_1,\dots,n_L$ enteros positivos. Un \emph{perceptrón multicapa} (MLP) de $L$ capas con anchuras $n_0\to n_1\to\dots\to n_L$ está formado por matrices $W_\ell\in\R^{n_{\ell-1}\times n_\ell}$, sesgos $b_\ell\in\R^{1\times n_\ell}$ y activaciones $\phi_\ell$, $\ell=1,\dots,L$. Para $x\in\R^{1\times n_0}$ se define $h^{(0)}=x$ y
\[
z^{(\ell)}=h^{(\ell-1)}W_\ell+b_\ell,\qquad h^{(\ell)}=\phi_\ell\bigl(z^{(\ell)}\bigr),\qquad \ell=1,\dots,L,
\]
con $\phi_\ell$ aplicada componente a componente. Las capas $1,\dots,L-1$ se llaman \emph{ocultas}. En un clasificador de $K$ clases, $n_L=K$, $\phi_L$ es la identidad y la salida $z^{(L)}$ es el vector de logits, que se convierte en $p=\softmax(z^{(L)})$.
\end{definicion}

\addtocounter{figure}{1}%

Con un lote de $B$ ejemplos apilados como filas de una matriz $X\in\R^{B\times n_0}$, la misma fórmula $XW_\ell+b_\ell$ (con $b_\ell$ sumado a cada fila) calcula todas las preactivaciones de la capa con un único producto de matrices. Esa es la razón práctica de la convención de vectores fila.

\begin{proposicion}[Número de parámetros]\label{prop:parametros-mlp}
Un MLP de anchuras $n_0\to n_1\to\dots\to n_L$ tiene $\sum_{\ell=1}^{L}(n_{\ell-1}+1)\,n_\ell$ parámetros.
\end{proposicion}

\addtocounter{definicion}{1}%

\begin{proposicion}[Colapso de las capas lineales]\label{prop:colapso}
Si todas las activaciones de un MLP son la identidad, existen $W\in\R^{n_0\times n_L}$ y $b\in\R^{1\times n_L}$ tales que $z^{(L)}=xW+b$ para toda entrada $x$. Además, $\operatorname{rg}W\le\min_{\ell}n_\ell$.
\end{proposicion}

La \cref{prop:colapso} explica por qué la no linealidad es imprescindible: sin ella, añadir capas no aumenta lo que la red puede representar, y una capa intermedia estrecha incluso lo reduce.

\addtocounter{definicion}{1}%

